
\begin{lemma}[Leibniz Rule]\label{lem:leibniz}
	Let $U \subset \mathbb{R}^{n}$ open, $\alpha  \in \Omega^{k}(U)$ and $\beta  \in \Omega^{\ell}(U)$. Then
	\begin{align*}
		d(\alpha \wedge \beta ) = (d\alpha) \wedge \beta + (-1)^{k} \alpha \wedge d\beta.
	\end{align*}
\end{lemma}
\begin{proof}
	Denote $\alpha = \sum_{I \in \mathcal{I}_{k}} a_{I} e_{I}^*$ and $\beta = \sum_{J \in \mathcal{I}_{\ell}} b_{J} e_{J}^*$. By definition,
	\begin{align*}
		d\alpha = \sum_{I \in \mathcal{I}_{k}} d a_{I} \wedge e_{I}^*, \quad d \beta = \sum_{J \in \mathcal{I}_{\ell}} d b_{J} \wedge e_{J}^*.
	\end{align*}
	On the other hand, by simple product rule $d(a_{I} b_{I}) = (da_{I}) b_{I} + a_{I} (db_{I})$,
	\begin{align*}
		d(\alpha \wedge \beta ) = \sum_{I \in \mathcal{I}_{k}} \sum_{J \in \mathcal{I}_{\ell}} (d a_{I} b_{J} + a_{I} db_{J}) \, e_{I}^* \wedge e_{J}^*.
	\end{align*}
	We can split this up into
	\begin{align*}
		\sum_{I \in \mathcal{I}_{k}} \sum_{J \in \mathcal{I}_{\ell}} d a_{I} b_{I} \, e_{I}^* \wedge e_{J}^* +
		\sum_{I \in \mathcal{I}_{k}} \sum_{J \in \mathcal{I}_{\ell}} a_{I} d b_{I} \, e_{I}^* \wedge e_{J}^*
	\end{align*}
	On the other hand,
	\begin{align*}
		d\alpha \wedge \beta + \alpha \wedge d\beta = \sum_{I \in \mathcal{I}_{k}} \sum_{J \in \mathcal{I}_{\ell}} da_{I} \,e_{I}^* \wedge (b_{J}\, e_{J}^*)  + \sum_{I \in \mathcal{I}_{k}} \sum_{J \in \mathcal{I}_{\ell}} a_{I}\, e_{I}^* \wedge (db_{J} \wedge e_{J}^*).
	\end{align*}
	Now, in the first sum, we can by definition of wedge commute $e_{I}^* \wedge (b_{J} \, e_{J}^*) = b_{J}\, e_{I}^* \wedge e_{J}^*$. Since the wedge is associative, we look for $e_{I}^* \wedge db_{J}$ and notice that
	\begin{align*}
		\det ((e_{i_{1}}, \dots, e_{i_{k}}, db_{J})^\top  A) = e_{I}^* \wedge db_{J}(A) = (-1)^{k} db_{J} \wedge e_{I}^*(A) = \det ((db_{I}, e_{i_{1}}, \dots, e_{i_{k}})^\top A),
	\end{align*}
	where we used the fact that the determinant is alternating in the columns of $(e_{i_{1}}, \dots, e_{i_{k}}, db_{J})$ (verify as an exercise). So we commute at the cost of the sign and get the statement.
\end{proof}

\begin{corollary}[Differential of wedged exact forms]\label{diff-exact-forms}
	Let $\alpha$ be a $k$-form such that $\alpha = d f_{1} \wedge d f_{2} \wedge \dots \wedge  df_{k}$ for $f_{i} \in C^{\infty}(U)$. Then $d\alpha = 0$.
\end{corollary}
\begin{proof}
	Apply Lemma \ref{lem:leibniz} and \ref{lem:double-diff} and prove by induction.
\end{proof}

\begin{definition}[Pullback]\label{def:pullback}
	Let $U \subset \mathbb{R}^{n}, V \subset \mathbb{R}^{m}$ open and $\alpha \in \Omega^{k}(U)$. Let $F: V \to U$ be a $C^{\infty}$ map. We define the pullback of $\alpha $ denoted $F^*\alpha \in \Omega^{k}(V)$ as follows:
	\begin{align*}
		(F^*\alpha)_{y}(A) := \alpha_{F(y)}(DF(y) \cdot A).
	\end{align*}
	It is an exercise to show that $F^*\alpha $ is indeed a $k$-form on $V$.
\end{definition}

\begin{lemma}[Pullback of basis form]\label{lem:pullback-basis}
	Let $U \subset \mathbb{R}^{n}$ open. Assume $\alpha = e_{I}^* \in \Omega^{k}(U)$ for $I = (i_{1}, \dots, i_{k}) \in \mathcal{I}_{k}$ is a $k$-form and $F(y) = (F_{1}(y), \dots, F_{n}(y))$ such that each $F_{i} \in C^{\infty}(U)$. Then
	\begin{align*}
		F^*\alpha = F^*e_{I}^* = dF_{i_{1}} \wedge \cdots \wedge dF_{i_{k}}.
	\end{align*}
	Sometimes we write $F^*dx_{I} = dF_{i_{1}} \wedge \cdots \wedge dF_{i_{k}}$.
\end{lemma}
\begin{proof}
	We have by definitions and computation of matrices,
	\begin{align*}
		(F^*\alpha)_{y} (A) & = \alpha_{F(y)}(DF(y) \cdot A) = e_{I}^*(DF(y) \cdot A) = \det (e_{I}^\top DF(y) \cdot A) \\ &= \det \left( (dF_{i_{1}}(y) \mid \cdots \mid dF_{i_{k}}(y))^\top \cdot A\right).
	\end{align*}
	Hence if you just split this into wedge factors, you get the statement.
\end{proof}

\begin{prop}[Pullback commutes with wedge and differential]\label{prop:pullback-commutes}
	Let $U \subset \mathbb{R}^{n}, V \subset \mathbb{R}^{m}$ open and $F: V \to U$ a $C^{\infty}$ map. Then
	\begin{enumerate}
		\item $\forall \alpha \in \Omega^{k}(U), \beta \in \Omega^{\ell}(U): F^*(\alpha \wedge \beta ) = F^*\alpha \wedge F^*\beta $,
		\item $\forall \alpha \in \Omega^{k}(U): d(F^*\alpha ) = F^*(d \alpha )$.
	\end{enumerate}
\end{prop}
\begin{proof}
	It once again is enough to show the statements for $\alpha_{x} = f(x) \,e_{I}^*$ and $\beta_{x} = g(x) \, e_{J}^*$ with $f,g: U \to \mathbb{R}$ $C^{\infty}$ functions and $I \in \mathcal{I}_{k}, J \in \mathcal{I}_{\ell}$.

	(1) By definition of wedge product, $(\alpha \wedge \beta )_{x} = f(x)g(x)\, e_{I}^* \wedge e_{J}^*$.
	By definition of pullback,
	\begin{align*}
		F^*(\alpha \wedge \beta)_y(T)
		= f(F(y))g(F(y)) \, (e_{I}^* \wedge e_{J}^*)(DF(y)\cdot T).
	\end{align*}
	On the other hand,
	\begin{align*}
		(F^*\alpha \wedge F^*\beta )_y(T)
		= f(F(y)) \, e_{I}^*(DF(y)\cdot T_{1}) \wedge g(F(y)) \, e_{J}^*(DF(y)\cdot T_{2}),
	\end{align*}
	if we split up $T \in \mathbb{R}^{n \times (k+\ell)}$ as $T = (T_{1}, T_{2})$ for $T_{1} \in \mathbb{R}^{n \times k}$ and $T_{2} \in \mathbb{R}^{n \times \ell}$. Using the wedge product definition on basis forms, we obtain
	\begin{align*}
		e_{I}^*(DF(y) T_{1}) \wedge e_{J}^*(DF(y) T_{2})
		 & =
		\det\!\begin{pmatrix}
			      e_{I}^\top DF(y)T_{1} \\
			      e_{J}^\top DF(y)T_{2}
		      \end{pmatrix}.
	\end{align*}
	Since $f(F(y))g(F(y))$ is scalar-valued and factors out, we conclude
	\begin{align*}
		(F^*\alpha \wedge F^*\beta )_y(T)
		 & = f(F(y))g(F(y)) \,
		\det\!\begin{pmatrix}
			      e_{I}^\top DF(y) T_{1} \\
			      e_{J}^\top DF(y) T_{2}
		      \end{pmatrix}
		\\ &
		=
		f(F(y)) g(F(y)) \det \! \left(
		\begin{pmatrix}
			e_{I}^\top \\
			e_{J}^\top
		\end{pmatrix}
		DF(y)T \right)
		F^*(\alpha \wedge \beta)_y(T).
	\end{align*}

	(2) We write $\alpha = f \, dx_{i_{1}} \wedge \cdots \wedge dx_{i_{k}}$. By definition of differential,
	$d\alpha = df \wedge dx_{i_{1}} \wedge \cdots \wedge dx_{i_{k}}$, hence by (1) and Lemma \ref{lem:pullback-basis},
	\begin{align*}
		F^*(d\alpha ) = F^*(df) \wedge F^*(dx_{i_{1}} \wedge \cdots \wedge dx_{i_{k}}) = F^*(df) \wedge dF_{i_{1}} \wedge \cdots \wedge dF_{i_{k}}.
	\end{align*}
	By the chain rule,
	\begin{align*}
		F^*(df)_{y}(T) = df_{F(y)}(DF(y) \cdot T) = Df_{F(y)} DF(x) \cdot T = D(f \circ F)_{y}T.
	\end{align*}
	Thereby $F^*(d\alpha ) = d(f \circ F) \wedge dF_{i_{1}} \wedge \cdots \wedge dF_{i_{k}}$. Meanwhile, splitting up the differential into its rows,
	\begin{align*}
		(F^*\alpha)_{y}(A) = \alpha_{F(y)}(DF(y) \cdot A) = (f \circ F \, dF_{i_{1}} \wedge \cdots \wedge dF_{i_{k}})_{y}(A).
	\end{align*}
	So we conclude $d(F^*\alpha ) = d(f \circ F) \wedge dF_{i_{1}} \wedge \cdots \wedge dF_{i_{k}}$ and hence $d(F^*\alpha ) = F^*(d \alpha )$.
\end{proof}

Now we bring this all back to the motivation and why we started inspecting differential forms in the first place. The key desideratum was invariance under reparametrisation. We however, consider only positive reparametrisations. What that means on a mathematical level is that we can drop the absolute value in the change of variables formula, and on a geometric level this imposes that we don't "change directions" in the middle of reparametrising. 

\begin{definition}[Orientable submanifold]\label{def:orientable-subm}
	A $k$-submanifold $M \subset \mathbb{R}^{n}$ is called orientable if it is covered by local parametrisations with transition maps that have positive Jacobi determinants. That is, for any two local parametrisations $\phi_{1}, \phi_{2}$ that cover different parts of the manifold $M \cap U_{1}$ and $M \cap U_{2}$ respectively, the reparametrisation $\psi$ from Remark \ref{rem:reparam} for $M \cap U_{1} \cap U_{2}$ always has positive Jacobi determinant $\det D\psi > 0$.

	An \textit{orientation} on a submanifold is then a fixed set of parametrisations $\phi_{i}$ that have positive transition maps. Any orientable submanifold is immediately assigned an orientation, unless specified otherwise. Once a set of parametrisations is chosen, the manifold is \textit{oriented}.
\end{definition}

\begin{definition}[Integral over parametrised submanifolds using $k$-forms]\label{def:int-param-subm-k-form}
	Let $V \subset \mathbb{R}^{k}$, $U \subset \mathbb{R}^{n}$ open and $\phi: V \to U$ be a parametrised submanifold $M = \phi(V)$ that is orientable. For $\alpha \in \Omega^{k}(U)$ with support in $V$, we define
	\begin{align*}
		\int_{M} \alpha := \int_{V} \phi^* \alpha =  \int_{V} \alpha_{\phi(y)} (D \phi(y)) \, dy,
	\end{align*}
	where $\phi^*\alpha = \alpha_{\phi(y)}(D \phi(y))$ is evaluated at the matrix $I_{n}$. This is invariant under positive reparametrisation by construction of differential forms (see beginning of section). 
\end{definition}

